\subsubsection{Instrumente und die genaue
Ancilla-Grenze}\label{instrumente-und-die-genaue-ancilla-grenze}

Ein Instrument mit Krausoperatoren Kx aus A0 kann aus einem scharfen
N-Sektor keinen anderen erzeugen. Eine spurtreue Summe solcher Zweige
erhält die gesamte Sektorverteilung. Selektive Konditionierung kann die
Gewichte einer bereits vorhandenen Mischung verändern, aber keine neue
Sektorstütze erzeugen.

Dasselbe gilt für einen zusätzlichen Zeiger, dessen Ladungsoperator auf
seinem ganzen Hilbertraum null ist, sofern auch die gemeinsame Kopplung
N erhält. Es gilt ausdrücklich \textbf{nicht} für eine beliebige
mitgeführte Ladungsreferenz. Bei deren Ladungswechsel kann die
Systemladung wechseln, während die Gesamtladung unverändert bleibt.
Abschnitt 5 gibt dafür einen expliziten positiven Zeugen. Ein Zeiger,
der anfangs Ladung null hat, aber andere Ladungszustände besitzt, ist
nicht mit einem auf dem ganzen Raum ladungstrivialen Zeiger
gleichzusetzen.

\subsection{3. Die kleinste zusätzliche Operation ist schon algebraisch
ausreichend}\label{die-kleinste-zusuxe4tzliche-operation-ist-schon-algebraisch-ausreichend}

In der ursprünglichen Basis ist

\[\adjustbox{max width=\linewidth}{$\displaystyle 
B_+=\frac12 b^\dagger\eta b^\dagger
=\sum_{A<\bar A}\eta_A b_A^\dagger b_{\bar A}^\dagger,
\qquad Q=B_++B_-.
$}\]

Dabei ist der zweite Bosonindex BAR(A). Es gibt dreißig solche
ungeordneten Paare. Die vollständige Invarianz von η wird an allen 45+15
ursprünglichen Liegeneratoren geprüft. Die 60 verschiedenen
Bosongewichte und der zusammenhängende Lie-Wirkungsgraph beweisen
zusätzlich, dass diese Paarung bis auf einen Skalar eindeutig ist.

Unter den G-invarianten, fermionparitätsgeraden, normalgeordneten
Polynomen ist Grad zwei der kleinste mögliche Grad einer N-ändernden
Operation: Es gibt kein lineares Bosonsingulett; fermionische Paarterme
haben nichttriviale SU(4)-Zentrumsladung; die verbleibenden Bilineare
f†f und b†b erhalten N. Die minimale Klasse ist daher κB+ + κ̄B−. Eine
Phase, Stärke oder Ausführungszeit wird dadurch nicht gewählt.

Die exakten CCR liefern am gesamten 60-Kanal-Tensor

\[\adjustbox{max width=\linewidth}{$\displaystyle 
[T_-,B_+]=R_+,\qquad
[Q,X]=R_--R_+,\qquad
\boxed{R_++R_-=-[N_b,[Q,X]].}
$}\]

Der ursprünglich zusätzlich vorgeschlagene kubische Quartettkanal
braucht somit keinen zweiten unabhängigen Koeffiziententensor. Ein
einziger zusätzlich gewährter Bosonpaargriff Q reicht algebraisch,
sobald X und Nb separat zugänglich sind. Der Nachweis benutzt exakte
Normalordnung auf dem unendlichen Boson-Fockraum, keine abgeschnittenen
Oszillator-CCR.

Das ist eine konkrete minimale \textbf{Erweiterung} des
Operationsvertrags. Die Quelle liefert η und W; sie liefert in der
geprüften Kette keine Ausführung von Q. Aus der Kommutatoridentität
allein folgt außerdem kein endliches exaktes Pulswort für exp(−itR).
Eine Interpretation als Kontrollsynthese benötigt schaltbare
Vorzeichen/Zeiten und eine bewiesene Konvergenz samt Fehlerkontrolle.
Bereits der statische Zusatz κQ bricht N; die separate
X/Nb-Verfügbarkeit ist nur für die genannte Synthese nötig.

\subsection{4. Ein neuer kleiner ursprünglicher
Singulettmechanismus}\label{ein-neuer-kleiner-urspruxfcnglicher-singulettmechanismus}

Definiere am wirklichen Vakuum der 64f/60b-Bank

\[\adjustbox{max width=\linewidth}{$\displaystyle 
|\beta\rangle=\frac{B_+|0\rangle}{\sqrt{30}},\qquad
|\rho\rangle=\frac{R_+|0\rangle}{\sqrt{480}}.
$}\]

Beide Zustände sind volle G-Singuletts und haben N=4. β enthält zwei
Bosonen; ρ enthält einen Boson und zwei Fermionen. Ihre Orthogonalität
folgt schon aus den unterschiedlichen Besetzungszahlen.

Der Schließungsschritt beruht auf der tatsächlichen Fierz-Identität

\[\adjustbox{max width=\linewidth}{$\displaystyle 
\sum_{AB}\eta_{AB}P_A^\dagger P_B^\dagger=0.
$}\]

Der Prüfer enumeriert hierfür \textbf{3840 signierte Quellterme auf 960
Fermionquartetts}. Nach vollständiger CAR-Normalordnung bleibt kein
einziger Koeffizient übrig. Zusätzlich wird die gesamte ursprüngliche
Wechselwirkung direkt auf alle 30 Komponenten von B+\textbar0⟩ und alle
480 Komponenten von R+\textbar0⟩ angewendet. Das ergibt genau

\[\adjustbox{max width=\linewidth}{$\displaystyle 
\begin{aligned}
T_+B_+|0\rangle&=0,&T_-B_+|0\rangle&=R_+|0\rangle,\\
T_+R_+|0\rangle&=16B_+|0\rangle,&T_-R_+|0\rangle&=0.
\end{aligned}
$}\]

Es gibt keine ausgelassenen Übergänge. Nach Normierung beträgt die
Kopplung in beiden Richtungen 4g. Deshalb ist span\{β,ρ\} unter dem
ganzen Hnat invariant und besitzt die eingangs angegebene
Zweizustandsmatrix. Eine Behauptung über die Dimension des gesamten
N=4-Singulettsektors ist für diesen Schluss nicht erforderlich und wird
hier nicht erhoben.

Für einen anfangs verfügbaren β-Zustand lautet die Umwandlung

\[\adjustbox{max width=\linewidth}{$\displaystyle 
P_{\beta\to\rho}(t)=
\frac{64g^2}{\Delta^2+64g^2}
\sin^2\!\left(\frac{t}{2}\sqrt{\Delta^2+64g^2}\right).
$}\]

Am ursprünglichen Prüfpunkt g/Δ=1/20 ist ihr Maximum exakt
\textbf{4/29}. Das ist Veränderung der Zusammensetzung in einer Bank;
die Rechnung führt keine Orte ein. Dieser N=4-Startzustand ist auch
nicht der frühere N=64-Grundzustand Ω.

Zum Vergleich bleibt der elementare N=2-Kanal W†\textbar A⟩/√8 ↔
\textbar bA⟩ mit Matrix {[}{[}0,√8g{]},{[}√8g,Δ{]}{]} erhalten. Sein
Maximum beträgt am selben Prüfpunkt 2/27; ein einzelnes ursprüngliches
Paar besitzt nur 1/8 des hellen Gewichts und erreicht 1/108. Auch hier
sind Vorbereitung, Auslesung und Zeitkalibrierung zusätzliche Teile
eines ausgeführten Prozessvertrags.

\subsection{5. Exakter gemeinsamer Dreizustandsprozess mit voller
innerer
Symmetrie}\label{exakter-gemeinsamer-dreizustandsprozess-mit-voller-innerer-symmetrie}

Man gewähre eine Referenz R mit zwei \textbf{G-trivialen} Zuständen
\textbar0⟩R und \textbar4⟩R, deren N\_R-Ladungen null beziehungsweise
vier sind. Setze

\[\adjustbox{max width=\linewidth}{$\displaystyle 
L_-=|0\rangle_R\langle4|,\qquad
V=\kappa(B_+\otimes L_-+B_-\otimes L_-^\dagger),
$}\]

\[\adjustbox{max width=\linewidth}{$\displaystyle 
H_{\rm joint}=H_{\rm nat}\otimes I+
E_R I\otimes|4\rangle_R\langle4|+V.
$}\]

Für reelle g, κ und ER ist Hjoint auf jedem festen Gesamt-N-Sektor eine
endliche hermitesche Matrix. V erhält Ngesamt und die \textbf{ganze}
innere Gruppe: B± sind bereits G-Skalare, die Referenz ist G-trivial.
Anders als ein ausschließlich auf Cartanladungen beruhender Ansatz ist
hier kein ungeprüfter Rest der nichtabelschen Gruppe übrig.

Der gemeinsame Raum

\[\adjustbox{max width=\linewidth}{$\displaystyle 
\mathcal K=\operatorname{span}\{
|0\rangle_S|4\rangle_R,
|\beta\rangle_S|0\rangle_R,
|\rho\rangle_S|0\rangle_R\}
$}\]

hat durchgehend Ngesamt=4 und ist unter Hjoint invariant. Denn
B−β=√30\textbar0⟩, B−ρ=0, L−\textbar0⟩R=0 und die nativen Übergänge sind
vollständig in Abschnitt 4 bewiesen. In dieser orthonormalen Basis
lautet der \textbf{exakte ganze Block}

\[\adjustbox{max width=\linewidth}{$\displaystyle 
\boxed{
H_{\rm joint}|_{\mathcal K}=
\begin{pmatrix}
E_R&\sqrt{30}\kappa&0\\
\sqrt{30}\kappa&2\Delta&4g\\
0&4g&\Delta
\end{pmatrix}.}
$}\]

Die zweite Kopplung ist unverändert aus Hnat abgeleitet. Die erste
stammt aus dem ausdrücklich neuen gemeinsamen Instrument V. Für den
Start \textbar0⟩S\textbar4⟩R ist die Quartettwahrscheinlichkeit bei
kleinen Zeiten

\[\adjustbox{max width=\linewidth}{$\displaystyle 
P_{\rho,0_R}(t)=120\kappa^2g^2t^4+O(t^6).
$}\]

Das positive führende Glied ist exakt aus (Hjoint²)31 berechnet; die
Geradheit folgt bei reellen Parametern aus der reellen symmetrischen
Matrix. Es wird weder eine besondere Auswahl der Energie ER noch ein
universeller hoher Übertragungsgrad behauptet.

\subsubsection{Die Referenz verschwindet nicht aus der
Rechnung}\label{die-referenz-verschwindet-nicht-aus-der-rechnung}

Ein allgemeiner gemeinsamer Zustand in K hat die Form

\[\adjustbox{max width=\linewidth}{$\displaystyle 
a|0,4_R\rangle+b|\beta,0_R\rangle+c|\rho,0_R\rangle.
$}\]

Nach Ausspuren der Referenz bleibt

\[\adjustbox{max width=\linewidth}{$\displaystyle 
\rho_S=|a|^2|0\rangle\langle0|+
(b|\beta\rangle+c|\rho\rangle)
(\bar b\langle\beta|+\bar c\langle\rho|).
$}\]

Die N=0/N=4-Kohärenzen sind null; die β/ρ-Kohärenz innerhalb N=4 bleibt.
Ein erfolgreicher Ladungseintrag hinterlässt die Referenz bei null statt
vier. Der Prozess implementiert deshalb keinen unveränderten
katalytischen Q-Antrieb des isolierten Systems. Eine zusätzliche
kohärente Referenz kann einen effektiven Antrieb in erster Ordnung
liefern; eine beliebige endliche Referenz ist aber nicht automatisch ein
exakter wiederverwendbarer Antrieb. Ein ergänzender exakter
Vierzustandszeuge zeigt beides: Eine scharfe Ladungsreferenz erzeugt
eine inkohärente Populationsänderung; für die gleichgewichtete kohärente
Referenz hat die reduzierte Systemdichte nach dem gewählten Puls
Determinante 1/16 und ist somit gemischt.

\subsubsection{Anschluss an das nachgereichte
Forschungspaper}\label{anschluss-an-das-nachgereichte-forschungspaper}

Abschnitt 9 des nachgereichten Papers benutzt dasselbe Prinzip einer
mitgerechneten Ladungsreferenz. Sein dortiger N=9-Vierzustandstransfer
erhält mit der neuen Referenz \textbf{eine einzelne innere
Cartanladung}. Unsere Konstruktion betrifft stattdessen N und erhält die
vollständige innere Gruppe, erzeugt aber keine räumliche Übertragung.
Sie ersetzt den dortigen Transferzeugen nicht und beweist dessen native
Herkunft nicht. Das Paper benennt seine Ein-Cartan-Grenze und den
Referenzverbrauch selbst ausdrücklich. Hier wurde nur dieser begrenzte
Anschluss gelesen; sein Gesamtaudit liegt außerhalb dieses Strangs.

\subsection{6. Ressourcenbilanz und der jetzt konkrete nächste
Nachweis}\label{ressourcenbilanz-und-der-jetzt-konkrete-nuxe4chste-nachweis}

{\def\LTcaptype{none} % do not increment counter
\begin{longtable}[]{@{}
  >{\raggedright\arraybackslash}p{(\linewidth - 4\tabcolsep) * \real{0.3333}}
  >{\raggedright\arraybackslash}p{(\linewidth - 4\tabcolsep) * \real{0.3333}}
  >{\raggedright\arraybackslash}p{(\linewidth - 4\tabcolsep) * \real{0.3333}}@{}}
\toprule\noalign{}
\begin{minipage}[b]{\linewidth}\raggedright
Ressource
\end{minipage} & \begin{minipage}[b]{\linewidth}\raggedright
Herkunft im geprüften Vertrag
\end{minipage} & \begin{minipage}[b]{\linewidth}\raggedright
Was noch zu liefern wäre
\end{minipage} \\
\midrule\noalign{}
\endhead
\bottomrule\noalign{}
\endlastfoot
W, η, innere Darstellung, endliche Clock & Tatsächliche
Quellkonstruktion und erneuter exakter Replay & Physische Identifikation
des Trägers \\
Hnat mit Δ,g & Bestehender ausdrücklich festgelegter Modellvertrag &
Physische Auswahl und Kalibrierung \\
Zweizustandsblock β↔ρ & Neue vollständige Anwendung aller nativen Terme
& Ursprüngliche Vorbereitung von β oder ρ \\
Unabhängige X/Nb-Kontrollen & Gewährte Kontrollstufe, nicht aus fixes H
abgeleitet & Instrumente und Schaltvertrag \\
Q oder der gemeinsame Übergang V & Hier genau ausgeschriebene
zusätzliche Operation & Compilerwort oder Instrument, das diesen
Übergang wirklich erzeugt \\
G-triviale Referenz mit N\_R=0,4 & Minimaler zusätzlicher
Ladungsspeicher für diese Konstruktion & Träger, Präparation, Energie ER
und alle Ressourcenkosten \\
Anfang & Vakuum × scharfe Referenzladung vier & Herkunft dieses
Gesamtzustands \\
Ausgang & Gemeinsame Besetzungen und Referenzwechsel & Ausführbare
gemeinsame Auslesung, ohne stilles Ausspuren relevanter Information \\
\end{longtable}
}

Die Minimalität ist präzise begrenzt: Grad zwei ist der kleinste
zugelassene G-invariante paritätsgerade N-ändernde Systemgrad; zwei
verschiedene Referenzladungen sind die kleinste Dimension für den
einmaligen scharfen Ausgleich; die Dreizustandsmatrix ist der kleinste
von diesem konkreten Anfang durch die beiden nichtverschwindenden
Kopplungen erzeugte Raum. Das ist keine globale Minimierung aller
möglichen Universen oder Compiler.

\textbf{Nächster Herkunftsnachweis:} Ein vorhandener ursprünglicher
Quellbaustein muss einen getypten Adapter zu V einschließlich
Referenzzustand und gemeinsamer Auslesung liefern. Ein Vorschlag allein
aus den bereits N-neutralen Fockwörtern scheitert jetzt am bewiesenen
Erhaltungssatz und braucht keine weitere Wortsuche. Ein anderer getypter
Quellbaustein bleibt eine offene Möglichkeit; bloße E8-Grade oder
gleiche Matrizenabmessungen reichen als Adapter nicht aus.

\subsection{7. Reproduktion und
Beweisumfang}\label{reproduktion-und-beweisumfang}

Aus diesem Ordner:

\begin{verbatim}
python3 -B -W error verify.py --output results_normal.json
python3 -OO -B -W error verify.py --output results_optimized.json
cmp results_normal.json results_optimized.json
\end{verbatim}

Benötigt werden Python, NumPy, SciPy und SymPy. Der Replay liest nur die
eingefrorenen Eingaben und braucht weder das ursprüngliche Repository
noch Netzwerkzugriff. \codeword{freeze.py} ist nur das Herkunftswerkzeug
für die erste Archivierung und gehört nicht zum wissenschaftlichen
Replay.

Es gibt \textbf{261 explizite exakte Prüfbedingungen pro Modus},
zusätzlich zu den separat ausgewiesenen 17 erhaltenen nativen
Konstruktorguards. Die vier Clockguards sind in den 261 enthalten.
Normaler und \codeword{-OO}-Lauf erzeugen byteidentische JSON-Dateien.
Die Bosonidentitäten nutzen die exakte Weyl- Normalordnung; die nativen
Zustandsaktionen nutzen alle CAR-Vorzeichen und Bosonmultiplikitäten.
Keine Prüfung benutzt ein wegoptimierbares \codeword{assert}, und
Warnungen werden als Fehler behandelt.

Die Bedingungen prüfen konkrete algebraische Tatsachen und Quellenpins;
die allgemeinen Induktions- und Verfügbarkeitsargumente stehen oben. Die
Zahl 261 ist keine Zahl unabhängiger Theoreme oder physischer
Realisierungen. Kein neuer Grundzustandssatz für die Erweiterung, keine
native Raumzeit, keine vollständige E8-Operatordarstellung und kein
T1-T8- Abschluss wird behauptet. Fremde Quellen, v1.6.7, zentrale Ledger
und Webseiten wurden nicht verändert.

\begin{center}\rule{0.5\linewidth}{0.5pt}\end{center}

\section{Teilbericht B: vollständige innere
Schattenkarte}\label{teilbericht-b-vollstuxe4ndige-innere-schattenkarte}

\section{Gemeinsame Schattenkarte auf dem ursprünglichen
W}\label{gemeinsame-schattenkarte-auf-dem-urspruxfcnglichen-w}

Stand: 15. September 2026. Eigenständiger Folgeschritt zu v1.6.7.

\subsection{Ergebnis}\label{ergebnis-2}

Der ursprüngliche Tensor liefert jetzt eine vollständig berechnete
gemeinsame Schattenkarte für den inneren hellen N=2-Code. Einzelne
Besetzungen unterscheiden nur 24, sämtliche Einteilchenoperatoren nur
736 von insgesamt 3600 hermiteschen Operatorrichtungen. Paarbesetzungen
zusammen mit ausführbaren Drehungen aus G=Spin(10)×SU(4) liefern dagegen
eine vollständige Zustandsrekonstruktion. Eine konkrete endliche Familie
mit 3599 nichtkonstanten Erwartungswerten liegt bei. Diese letzte
Aussage ist \textbf{bedingt auf die Ausführbarkeit der zusätzlichen
Kontrollen und Messungen}. Aus ihrer algebraischen Existenz folgt noch
kein ursprüngliches Messinstrument samt Zustandsänderung.

Das Ergebnis betrifft denselben ursprünglichen W, keine ersatzweise
eingeführte Qubit-Analogie. Der W-Dateipin ist
\codeword{3f00a0892157a691e4ef26920c3702772c7bf76a250fad98ff04a079bb64b763}.
Die vier gelesenen v1.6.7-Kontextdateien sind ebenfalls unverändert
eingefroren.

